\chapter{Unit 4: Array Operations using NumPy}

\section{Section A: Multiple Choice Practice Questions (60 Curated Problems)}

\noindent This section contains 60 rigorous multiple-choice questions extracted from university course examination assessments. Review each question, test edge-case syntax, and verify code execution mechanics.
\vspace{3mm}

\noindent\textbf{Q1.} \hfill {\small\color{LPUOrange}\textbf{[arrays vs lists | Easy]}}\\[2pt]
Which Python library provides the \texttt{ndarray} object used for numerical arrays?

\begin{enumerate}[label=(\alph*),topsep=2pt,itemsep=1pt,partopsep=0pt,parsep=0pt]
    \item \relax Matplotlib
    \item \relax NumPy
    \item \relax SciPy
    \item \relax Pandas
\end{enumerate}
\vspace{2mm}

\noindent\textbf{Q2.} \hfill {\small\color{LPUOrange}\textbf{[arrays vs lists | Easy]}}\\[2pt]
What is a key difference between a NumPy array and a Python list?

\begin{enumerate}[label=(\alph*),topsep=2pt,itemsep=1pt,partopsep=0pt,parsep=0pt]
    \item \relax A NumPy array cannot be indexed
    \item \relax A list uses less memory than an array
    \item \relax A NumPy array stores elements of a single data type
    \item \relax A list can only hold numbers
\end{enumerate}
\vspace{2mm}

\noindent\textbf{Q3.} \hfill {\small\color{LPUOrange}\textbf{[arrays vs lists | Easy]}}\\[2pt]
Which function is used to create a NumPy array from a Python list?

\begin{enumerate}[label=(\alph*),topsep=2pt,itemsep=1pt,partopsep=0pt,parsep=0pt]
    \item \relax \texttt{np.list()}
    \item \relax \texttt{np.make()}
    \item \relax \texttt{np.create()}
    \item \relax \texttt{np.array()}
\end{enumerate}
\vspace{2mm}

\noindent\textbf{Q4.} \hfill {\small\color{LPUOrange}\textbf{[arrays vs lists | Easy]}}\\[2pt]
Compared to Python lists, NumPy arrays are generally preferred for numerical computation because they are:

\begin{enumerate}[label=(\alph*),topsep=2pt,itemsep=1pt,partopsep=0pt,parsep=0pt]
    \item \relax Limited to one element only
    \item \relax Faster and more memory efficient
    \item \relax Unable to store integers
    \item \relax Slower but easier to read
\end{enumerate}
\vspace{2mm}

\noindent\textbf{Q5.} \hfill {\small\color{LPUOrange}\textbf{[data types | Easy]}}\\[2pt]
Which attribute is used to check the data type of elements in a NumPy array?

\begin{enumerate}[label=(\alph*),topsep=2pt,itemsep=1pt,partopsep=0pt,parsep=0pt]
    \item \relax \texttt{.datatype}
    \item \relax \texttt{.dtype}
    \item \relax \texttt{.type}
    \item \relax \texttt{.kind}
\end{enumerate}
\vspace{2mm}

\noindent\textbf{Q6.} \hfill {\small\color{LPUOrange}\textbf{[data types | Easy]}}\\[2pt]
What does the NumPy data type \texttt{int64} represent?

\begin{enumerate}[label=(\alph*),topsep=2pt,itemsep=1pt,partopsep=0pt,parsep=0pt]
    \item \relax A 64-character string
    \item \relax A 64-element array
    \item \relax A 64-bit floating point number
    \item \relax A 64-bit integer
\end{enumerate}
\vspace{2mm}

\noindent\textbf{Q7.} \hfill {\small\color{LPUOrange}\textbf{[data types | Easy]}}\\[2pt]
Which method changes the data type of an existing NumPy array?

\begin{enumerate}[label=(\alph*),topsep=2pt,itemsep=1pt,partopsep=0pt,parsep=0pt]
    \item \relax \texttt{.astype()}
    \item \relax \texttt{.convert()}
    \item \relax \texttt{.retype()}
    \item \relax \texttt{.setdtype()}
\end{enumerate}
\vspace{2mm}

\noindent\textbf{Q8.} \hfill {\small\color{LPUOrange}\textbf{[data types | Easy]}}\\[2pt]
If an array is created as \texttt{np.array([1.0, 2.0, 3.0])}, its default data type will be:

\begin{enumerate}[label=(\alph*),topsep=2pt,itemsep=1pt,partopsep=0pt,parsep=0pt]
    \item \relax \texttt{bool}
    \item \relax \texttt{str}
    \item \relax \texttt{float64}
    \item \relax \texttt{int32}
\end{enumerate}
\vspace{2mm}

\noindent\textbf{Q9.} \hfill {\small\color{LPUOrange}\textbf{[array operations | Easy]}}\\[2pt]
Given \texttt{a = np.array([1, 2, 3])}, what is the result of \texttt{a + 2}?

\begin{enumerate}[label=(\alph*),topsep=2pt,itemsep=1pt,partopsep=0pt,parsep=0pt]
    \item \relax \texttt{[3, 4, 5]}
    \item \relax \texttt{[1, 2, 3, 2]}
    \item \relax \texttt{[2, 4, 6]}
    \item \relax \texttt{[3, 2, 3]}
\end{enumerate}
\vspace{2mm}

\noindent\textbf{Q10.} \hfill {\small\color{LPUOrange}\textbf{[array operations | Easy]}}\\[2pt]
Which attribute returns the shape (dimensions) of a NumPy array?

\begin{enumerate}[label=(\alph*),topsep=2pt,itemsep=1pt,partopsep=0pt,parsep=0pt]
    \item \relax \texttt{.size}
    \item \relax \texttt{.shape}
    \item \relax \texttt{.dim}
    \item \relax \texttt{.length}
\end{enumerate}
\vspace{2mm}

\noindent\textbf{Q11.} \hfill {\small\color{LPUOrange}\textbf{[array operations | Easy]}}\\[2pt]
Given \texttt{a = np.array([1, 2, 3])} and \texttt{b = np.array([4, 5, 6])}, what is \texttt{a * b}?

\begin{enumerate}[label=(\alph*),topsep=2pt,itemsep=1pt,partopsep=0pt,parsep=0pt]
    \item \relax \texttt{[4, 5, 6]}
    \item \relax \texttt{[1, 2, 3]}
    \item \relax \texttt{[4, 10, 18]}
    \item \relax \texttt{[5, 7, 9]}
\end{enumerate}
\vspace{2mm}

\noindent\textbf{Q12.} \hfill {\small\color{LPUOrange}\textbf{[array operations | Easy]}}\\[2pt]
Which function reshapes an array without changing its data?

\begin{enumerate}[label=(\alph*),topsep=2pt,itemsep=1pt,partopsep=0pt,parsep=0pt]
    \item \relax \texttt{restack()}
    \item \relax \texttt{resize()}
    \item \relax \texttt{reform()}
    \item \relax \texttt{reshape()}
\end{enumerate}
\vspace{2mm}

\noindent\textbf{Q13.} \hfill {\small\color{LPUOrange}\textbf{[array operations | Easy]}}\\[2pt]
What does \texttt{np.arange(0, 5)} produce?

\begin{enumerate}[label=(\alph*),topsep=2pt,itemsep=1pt,partopsep=0pt,parsep=0pt]
    \item \relax \texttt{[0, 1, 2, 3, 4]}
    \item \relax \texttt{[0, 1, 2, 3, 4, 5]}
    \item \relax \texttt{[1, 2, 3, 4, 5]}
    \item \relax \texttt{[0, 5]}
\end{enumerate}
\vspace{2mm}

\noindent\textbf{Q14.} \hfill {\small\color{LPUOrange}\textbf{[statistical functions | Easy]}}\\[2pt]
Which NumPy function computes the average of array elements?

\begin{enumerate}[label=(\alph*),topsep=2pt,itemsep=1pt,partopsep=0pt,parsep=0pt]
    \item \relax \texttt{np.mean()}
    \item \relax \texttt{np.middle()}
    \item \relax \texttt{np.center()}
    \item \relax \texttt{np.avg()}
\end{enumerate}
\vspace{2mm}

\noindent\textbf{Q15.} \hfill {\small\color{LPUOrange}\textbf{[statistical functions | Easy]}}\\[2pt]
What does \texttt{np.max()} return for a NumPy array?

\begin{enumerate}[label=(\alph*),topsep=2pt,itemsep=1pt,partopsep=0pt,parsep=0pt]
    \item \relax The number of elements
    \item \relax The smallest element
    \item \relax The sum of elements
    \item \relax The largest element
\end{enumerate}
\vspace{2mm}

\noindent\textbf{Q16.} \hfill {\small\color{LPUOrange}\textbf{[statistical functions | Easy]}}\\[2pt]
Which function computes the standard deviation of a NumPy array?

\begin{enumerate}[label=(\alph*),topsep=2pt,itemsep=1pt,partopsep=0pt,parsep=0pt]
    \item \relax \texttt{np.dev()}
    \item \relax \texttt{np.std()}
    \item \relax \texttt{np.var()}
    \item \relax \texttt{np.spread()}
\end{enumerate}
\vspace{2mm}

\noindent\textbf{Q17.} \hfill {\small\color{LPUOrange}\textbf{[statistical functions | Easy]}}\\[2pt]
Given \texttt{a = np.array([2, 4, 6])}, what does \texttt{np.sum(a)} return?

\begin{enumerate}[label=(\alph*),topsep=2pt,itemsep=1pt,partopsep=0pt,parsep=0pt]
    \item \relax \texttt{3}
    \item \relax \texttt{12}
    \item \relax \texttt{6}
    \item \relax \texttt{4}
\end{enumerate}
\vspace{2mm}

\noindent\textbf{Q18.} \hfill {\small\color{LPUOrange}\textbf{[broadcasting | Easy]}}\\[2pt]
What does broadcasting in NumPy allow you to do?

\begin{enumerate}[label=(\alph*),topsep=2pt,itemsep=1pt,partopsep=0pt,parsep=0pt]
    \item \relax Sort arrays automatically
    \item \relax Print arrays to screen
    \item \relax Send data over a network
    \item \relax Operate on arrays of different shapes
\end{enumerate}
\vspace{2mm}

\noindent\textbf{Q19.} \hfill {\small\color{LPUOrange}\textbf{[broadcasting | Easy]}}\\[2pt]
When you add a scalar to a NumPy array, broadcasting causes the scalar to be:

\begin{enumerate}[label=(\alph*),topsep=2pt,itemsep=1pt,partopsep=0pt,parsep=0pt]
    \item \relax Ignored entirely
    \item \relax Applied to every element
    \item \relax Added only to the first element
    \item \relax Converted to a string
\end{enumerate}
\vspace{2mm}

\noindent\textbf{Q20.} \hfill {\small\color{LPUOrange}\textbf{[broadcasting | Easy]}}\\[2pt]
For broadcasting to work between two arrays, their dimensions must be:

\begin{enumerate}[label=(\alph*),topsep=2pt,itemsep=1pt,partopsep=0pt,parsep=0pt]
    \item \relax Prime numbers
    \item \relax Compatible or equal to 1
    \item \relax Both greater than 10
    \item \relax Exactly identical always
\end{enumerate}
\vspace{2mm}

\noindent\textbf{Q21.} \hfill {\small\color{LPUOrange}\textbf{[arrays vs lists | Medium]}}\\[2pt]
Given \texttt{import numpy as np}, what is the output of the following code? 
 python 
 lst = [1, 2, 3] 
 arr = np.array([1, 2, 3]) 
 print(lst * 2, arr * 2)

\begin{enumerate}[label=(\alph*),topsep=2pt,itemsep=1pt,partopsep=0pt,parsep=0pt]
    \item \relax \texttt{[2, 4, 6] [2 4 6]}
    \item \relax \texttt{[2, 4, 6] [1 2 3 1 2 3]}
    \item \relax \texttt{[1, 2, 3, 1, 2, 3] [1 2 3 1 2 3]}
    \item \relax \texttt{[1, 2, 3, 1, 2, 3] [2 4 6]}
\end{enumerate}
\vspace{2mm}

\noindent\textbf{Q22.} \hfill {\small\color{LPUOrange}\textbf{[arrays vs lists | Medium]}}\\[2pt]
Which statement best explains why NumPy arrays are generally faster than Python lists for numerical computation?

\begin{enumerate}[label=(\alph*),topsep=2pt,itemsep=1pt,partopsep=0pt,parsep=0pt]
    \item \relax Arrays are interpreted line by line while lists are compiled
    \item \relax Arrays store each element as a separate Python object with pointers
    \item \relax Arrays automatically use multiple CPU cores for every operation
    \item \relax Arrays store homogeneous data in contiguous memory and use vectorized C operations
\end{enumerate}
\vspace{2mm}

\noindent\textbf{Q23.} \hfill {\small\color{LPUOrange}\textbf{[data types | Medium]}}\\[2pt]
What will be the \texttt{dtype} of the array created below? 
 python 
 arr = np.array([1, 2, 3.0, 4])

\begin{enumerate}[label=(\alph*),topsep=2pt,itemsep=1pt,partopsep=0pt,parsep=0pt]
    \item \relax \texttt{float64}
    \item \relax \texttt{int64}
    \item \relax \texttt{object}
    \item \relax \texttt{complex128}
\end{enumerate}
\vspace{2mm}

\noindent\textbf{Q24.} \hfill {\small\color{LPUOrange}\textbf{[data types | Medium]}}\\[2pt]
What is the output of the following code? 
 python 
 arr = np.array([1, 2, 3], dtype=np.int8) 
 arr[0] = 300 
 print(arr[0])

\begin{enumerate}[label=(\alph*),topsep=2pt,itemsep=1pt,partopsep=0pt,parsep=0pt]
    \item \relax Raises an OverflowError
    \item \relax \texttt{300}
    \item \relax \texttt{127}
    \item \relax \texttt{44}
\end{enumerate}
\vspace{2mm}

\noindent\textbf{Q25.} \hfill {\small\color{LPUOrange}\textbf{[data types | Medium]}}\\[2pt]
Which code correctly converts a float array \texttt{a} into an integer array without changing the original?

\begin{enumerate}[label=(\alph*),topsep=2pt,itemsep=1pt,partopsep=0pt,parsep=0pt]
    \item \relax \texttt{b = a.dtype(np.int32)}
    \item \relax \texttt{b = a.convert('int32')}
    \item \relax \texttt{b = int(a)}
    \item \relax \texttt{b = a.astype(np.int32)}
\end{enumerate}
\vspace{2mm}

\noindent\textbf{Q26.} \hfill {\small\color{LPUOrange}\textbf{[array operations | Medium]}}\\[2pt]
Given \texttt{arr = np.arange(12).reshape(3, 4)}, what does \texttt{arr[:, 1]} return?

\begin{enumerate}[label=(\alph*),topsep=2pt,itemsep=1pt,partopsep=0pt,parsep=0pt]
    \item \relax \texttt{[1 2 3 4]}
    \item \relax \texttt{[1 5 9]}
    \item \relax \texttt{[0 1 2 3]}
    \item \relax \texttt{[4 5 6 7]}
\end{enumerate}
\vspace{2mm}

\noindent\textbf{Q27.} \hfill {\small\color{LPUOrange}\textbf{[array operations | Medium]}}\\[2pt]
What is the result of the following code? 
 python 
 a = np.array([[1, 2], [3, 4]]) 
 print(a.sum(axis=0))

\begin{enumerate}[label=(\alph*),topsep=2pt,itemsep=1pt,partopsep=0pt,parsep=0pt]
    \item \relax \texttt{[4 6]}
    \item \relax \texttt{[3 7]}
    \item \relax \texttt{[6 4]}
    \item \relax \texttt{10}
\end{enumerate}
\vspace{2mm}

\noindent\textbf{Q28.} \hfill {\small\color{LPUOrange}\textbf{[array operations | Medium]}}\\[2pt]
Given \texttt{arr = np.array([10, 20, 30, 40, 50])}, what does \texttt{arr[::-2]} return?

\begin{enumerate}[label=(\alph*),topsep=2pt,itemsep=1pt,partopsep=0pt,parsep=0pt]
    \item \relax \texttt{[10 30 50]}
    \item \relax \texttt{[40 20]}
    \item \relax \texttt{[50 30 10]}
    \item \relax \texttt{[50 40 30 20 10]}
\end{enumerate}
\vspace{2mm}

\noindent\textbf{Q29.} \hfill {\small\color{LPUOrange}\textbf{[array operations | Medium]}}\\[2pt]
What is the difference between \texttt{np.dot(a, b)} and \texttt{a * b} for two 2D arrays of the same shape?

\begin{enumerate}[label=(\alph*),topsep=2pt,itemsep=1pt,partopsep=0pt,parsep=0pt]
    \item \relax \texttt{np.dot} performs element-wise multiplication while \texttt{*} performs matrix multiplication
    \item \relax Both perform matrix multiplication but \texttt{*} is faster
    \item \relax Both perform element-wise multiplication with different rounding
    \item \relax \texttt{np.dot} performs matrix multiplication while \texttt{*} performs element-wise multiplication
\end{enumerate}
\vspace{2mm}

\noindent\textbf{Q30.} \hfill {\small\color{LPUOrange}\textbf{[array operations | Medium]}}\\[2pt]
What is the output of the following code? 
 python 
 arr = np.array([1, 2, 3, 4, 5]) 
 print(arr[arr \% 2 == 0])

\begin{enumerate}[label=(\alph*),topsep=2pt,itemsep=1pt,partopsep=0pt,parsep=0pt]
    \item \relax \texttt{[True False True False True]}
    \item \relax \texttt{[1 3 5]}
    \item \relax \texttt{[2 4]}
    \item \relax \texttt{[0 2 0 4 0]}
\end{enumerate}
\vspace{2mm}

\noindent\textbf{Q31.} \hfill {\small\color{LPUOrange}\textbf{[statistical functions | Medium]}}\\[2pt]
Given \texttt{arr = np.array([2, 4, 6, 8])}, what does \texttt{np.std(arr)} compute and its approximate value?

\begin{enumerate}[label=(\alph*),topsep=2pt,itemsep=1pt,partopsep=0pt,parsep=0pt]
    \item \relax The variance, 5.0
    \item \relax Sample standard deviation, $\approx$
    \item \relax Population standard deviation, $\approx$
    \item \relax The range, 6.0
\end{enumerate}
\vspace{2mm}

\noindent\textbf{Q32.} \hfill {\small\color{LPUOrange}\textbf{[statistical functions | Medium]}}\\[2pt]
For \texttt{a = np.array([[1, 2], [3, 4]])}, what does \texttt{np.mean(a, axis=1)} return?

\begin{enumerate}[label=(\alph*),topsep=2pt,itemsep=1pt,partopsep=0pt,parsep=0pt]
    \item \relax \texttt{[1.5 3.5]}
    \item \relax \texttt{2.5}
    \item \relax \texttt{[2. 3.]}
    \item \relax \texttt{[4. 6.]}
\end{enumerate}
\vspace{2mm}

\noindent\textbf{Q33.} \hfill {\small\color{LPUOrange}\textbf{[statistical functions | Medium]}}\\[2pt]
Which function returns the *index* of the maximum value in an array rather than the value itself?

\begin{enumerate}[label=(\alph*),topsep=2pt,itemsep=1pt,partopsep=0pt,parsep=0pt]
    \item \relax \texttt{np.maximum}
    \item \relax \texttt{np.argmax}
    \item \relax \texttt{np.max}
    \item \relax \texttt{np.amax}
\end{enumerate}
\vspace{2mm}

\noindent\textbf{Q34.} \hfill {\small\color{LPUOrange}\textbf{[statistical functions | Medium]}}\\[2pt]
What is the output of the following code? 
 python 
 arr = np.array([3, 1, 4, 1, 5, 9, 2]) 
 print(np.median(arr))

\begin{enumerate}[label=(\alph*),topsep=2pt,itemsep=1pt,partopsep=0pt,parsep=0pt]
    \item \relax \texttt{4.0}
    \item \relax \texttt{3.5}
    \item \relax \texttt{1.0}
    \item \relax \texttt{3.0}
\end{enumerate}
\vspace{2mm}

\noindent\textbf{Q35.} \hfill {\small\color{LPUOrange}\textbf{[statistical functions | Medium]}}\\[2pt]
Given \texttt{arr = np.array([10, 20, 30, 40])}, what does \texttt{np.cumsum(arr)} return?

\begin{enumerate}[label=(\alph*),topsep=2pt,itemsep=1pt,partopsep=0pt,parsep=0pt]
    \item \relax \texttt{[10 20 30 40]}
    \item \relax \texttt{100}
    \item \relax \texttt{[10 30 60 100]}
    \item \relax \texttt{[100 90 70 40]}
\end{enumerate}
\vspace{2mm}

\noindent\textbf{Q36.} \hfill {\small\color{LPUOrange}\textbf{[broadcasting | Medium]}}\\[2pt]
What is the result of the following code? 
 python 
 a = np.array([[1], [2], [3]]) 
 b = np.array([10, 20, 30]) 
 print((a + b).shape)

\begin{enumerate}[label=(\alph*),topsep=2pt,itemsep=1pt,partopsep=0pt,parsep=0pt]
    \item \relax \texttt{(3, 3)}
    \item \relax Raises a ValueError
    \item \relax \texttt{(3,)}
    \item \relax \texttt{(3, 1)}
\end{enumerate}
\vspace{2mm}

\noindent\textbf{Q37.} \hfill {\small\color{LPUOrange}\textbf{[broadcasting | Medium]}}\\[2pt]
Which pair of array shapes CANNOT be broadcast together?

\begin{enumerate}[label=(\alph*),topsep=2pt,itemsep=1pt,partopsep=0pt,parsep=0pt]
    \item \relax \texttt{(5, 1)} and \texttt{(1, 6)}
    \item \relax \texttt{(3, 4)} and \texttt{(4,)}
    \item \relax \texttt{(4, 3)} and \texttt{(2,)}
    \item \relax \texttt{(2, 3)} and \texttt{(1, 3)}
\end{enumerate}
\vspace{2mm}

\noindent\textbf{Q38.} \hfill {\small\color{LPUOrange}\textbf{[broadcasting | Medium]}}\\[2pt]
What is the output of the following code? 
 python 
 a = np.array([1, 2, 3]) 
 print(a + 5)

\begin{enumerate}[label=(\alph*),topsep=2pt,itemsep=1pt,partopsep=0pt,parsep=0pt]
    \item \relax \texttt{[1 2 3 5]}
    \item \relax \texttt{[5 5 5]}
    \item \relax \texttt{[6 7 8]}
    \item \relax Raises a ValueError
\end{enumerate}
\vspace{2mm}

\noindent\textbf{Q39.} \hfill {\small\color{LPUOrange}\textbf{[broadcasting | Medium]}}\\[2pt]
Given \texttt{a} of shape \texttt{(2, 3)} and \texttt{b} of shape \texttt{(3,)}, what is the shape of \texttt{a * b}?

\begin{enumerate}[label=(\alph*),topsep=2pt,itemsep=1pt,partopsep=0pt,parsep=0pt]
    \item \relax Raises a ValueError
    \item \relax \texttt{(2,)}
    \item \relax \texttt{(3,)}
    \item \relax \texttt{(2, 3)}
\end{enumerate}
\vspace{2mm}

\noindent\textbf{Q40.} \hfill {\small\color{LPUOrange}\textbf{[array operations | Medium]}}\\[2pt]
What is the output of the following code? 
 python 
 arr = np.array([1, 2, 3]) 
 arr2 = arr 
 arr2[0] = 99 
 print(arr[0])

\begin{enumerate}[label=(\alph*),topsep=2pt,itemsep=1pt,partopsep=0pt,parsep=0pt]
    \item \relax \texttt{0}
    \item \relax Raises an error
    \item \relax \texttt{1}
    \item \relax \texttt{99}
\end{enumerate}
\vspace{2mm}

\noindent\textbf{Q41.} \hfill {\small\color{LPUOrange}\textbf{[arrays vs lists | Hard]}}\\[2pt]
Consider the following code: 

\begin{pycode}
import numpy as np\\
lst = [1, 2, 3]\\
arr = np.array([1, 2, 3])\\
print(lst * 2)\\
print(arr * 2)
\end{pycode}

 What is the output?

\begin{enumerate}[label=(\alph*),topsep=2pt,itemsep=1pt,partopsep=0pt,parsep=0pt]
    \item \relax \texttt{[1, 2, 3, 1, 2, 3]} and \texttt{[1 2 3 1 2 3]}
    \item \relax \texttt{[2, 4, 6]} and \texttt{[1 2 3 1 2 3]}
    \item \relax \texttt{[1, 2, 3, 1, 2, 3]} and \texttt{[2 4 6]}
    \item \relax \texttt{[2, 4, 6]} and \texttt{[2 4 6]}
\end{enumerate}
\vspace{2mm}

\noindent\textbf{Q42.} \hfill {\small\color{LPUOrange}\textbf{[arrays vs lists | Hard]}}\\[2pt]
A NumPy array stores 1 0 6 integers while an equivalent Python list stores the same values. Why does the NumPy array typically consume far less memory?

\begin{enumerate}[label=(\alph*),topsep=2pt,itemsep=1pt,partopsep=0pt,parsep=0pt]
    \item \relax It stores only unique values and references duplicates
    \item \relax It offloads storage to disk and keeps only pointers in RAM
    \item \relax It compresses the integers using run-length encoding automatically
    \item \relax It stores raw values in a contiguous, fixed-type buffer without per-element object overhead
\end{enumerate}
\vspace{2mm}

\noindent\textbf{Q43.} \hfill {\small\color{LPUOrange}\textbf{[data types | Hard]}}\\[2pt]
What does the following code print? 

\begin{pycode}
import numpy as np\\
arr = np.array([1, 2, 3], dtype=np.int8)\\
arr[0] = 130\\
print(arr[0])
\end{pycode}

\begin{enumerate}[label=(\alph*),topsep=2pt,itemsep=1pt,partopsep=0pt,parsep=0pt]
    \item \relax \texttt{-126}
    \item \relax \texttt{130}
    \item \relax \texttt{127}
    \item \relax \texttt{2}
\end{enumerate}
\vspace{2mm}

\noindent\textbf{Q44.} \hfill {\small\color{LPUOrange}\textbf{[data types | Hard]}}\\[2pt]
Given \texttt{arr = np.array([1, 2, 3])}, which operation returns a new array of type \texttt{float64} without modifying \texttt{arr}?

\begin{enumerate}[label=(\alph*),topsep=2pt,itemsep=1pt,partopsep=0pt,parsep=0pt]
    \item \relax \texttt{arr.view(np.float64)}
    \item \relax \texttt{arr.astype(np.float64)}
    \item \relax \texttt{arr.dtype = np.float64}
    \item \relax \texttt{arr.reshape(np.float64)}
\end{enumerate}
\vspace{2mm}

\noindent\textbf{Q45.} \hfill {\small\color{LPUOrange}\textbf{[data types | Hard]}}\\[2pt]
What is the resulting dtype of \texttt{np.array([True, 2, 3.5])}?

\begin{enumerate}[label=(\alph*),topsep=2pt,itemsep=1pt,partopsep=0pt,parsep=0pt]
    \item \relax \texttt{int64}
    \item \relax \texttt{float64}
    \item \relax \texttt{object}
    \item \relax \texttt{bool}
\end{enumerate}
\vspace{2mm}

\noindent\textbf{Q46.} \hfill {\small\color{LPUOrange}\textbf{[array operations | Hard]}}\\[2pt]
What does the following print? 

\begin{pycode}
import numpy as np\\
a = np.arange(6).reshape(2, 3)\\
b = a.T\\
b[0, 0] = 99\\
print(a[0, 0])
\end{pycode}

\begin{enumerate}[label=(\alph*),topsep=2pt,itemsep=1pt,partopsep=0pt,parsep=0pt]
    \item \relax \texttt{0}
    \item \relax \texttt{99}
    \item \relax \texttt{5}
    \item \relax Raises an error
\end{enumerate}
\vspace{2mm}

\noindent\textbf{Q47.} \hfill {\small\color{LPUOrange}\textbf{[array operations | Hard]}}\\[2pt]
Given \texttt{a = np.array([10, 20, 30, 40])}, what does \texttt{a[::-1][1:3]} return?

\begin{enumerate}[label=(\alph*),topsep=2pt,itemsep=1pt,partopsep=0pt,parsep=0pt]
    \item \relax \texttt{array([30, 20])}
    \item \relax \texttt{array([20, 30])}
    \item \relax \texttt{array([40, 30])}
    \item \relax \texttt{array([30, 40])}
\end{enumerate}
\vspace{2mm}

\noindent\textbf{Q48.} \hfill {\small\color{LPUOrange}\textbf{[array operations | Hard]}}\\[2pt]
Consider: 

\begin{pycode}
import numpy as np\\
a = np.array([1, 2, 3, 4, 5])\\
mask = a \% 2 == 0\\
a[mask] = 0\\
print(a)
\end{pycode}

 What is the output?

\begin{enumerate}[label=(\alph*),topsep=2pt,itemsep=1pt,partopsep=0pt,parsep=0pt]
    \item \relax \texttt{[0 2 0 4 0]}
    \item \relax \texttt{[0 0 0 0 0]}
    \item \relax \texttt{[1 2 3 4 5]}
    \item \relax \texttt{[1 0 3 0 5]}
\end{enumerate}
\vspace{2mm}

\noindent\textbf{Q49.} \hfill {\small\color{LPUOrange}\textbf{[array operations | Hard]}}\\[2pt]
What is the result of \texttt{np.arange(12).reshape(3, 4)[:, 1::2]}?

\begin{enumerate}[label=(\alph*),topsep=2pt,itemsep=1pt,partopsep=0pt,parsep=0pt]
    \item \relax \texttt{[[2 3], [6 7], [10 11]]}
    \item \relax \texttt{[[1 3], [5 7], [9 11]]}
    \item \relax \texttt{[[0 2], [4 6], [8 10]]}
    \item \relax \texttt{[[1 2], [5 6], [9 10]]}
\end{enumerate}
\vspace{2mm}

\noindent\textbf{Q50.} \hfill {\small\color{LPUOrange}\textbf{[array operations | Hard]}}\\[2pt]
Which statement about \texttt{np.reshape(a, (-1, 3))} is correct when \texttt{a} has 12 elements?

\begin{enumerate}[label=(\alph*),topsep=2pt,itemsep=1pt,partopsep=0pt,parsep=0pt]
    \item \relax It raises an error because \texttt{-1} is not a valid dimension
    \item \relax The \texttt{-1} is treated as 0
    \item \relax The \texttt{-1} is inferred as 4
    \item \relax The \texttt{-1} reverses the array before reshaping
\end{enumerate}
\vspace{2mm}

\noindent\textbf{Q51.} \hfill {\small\color{LPUOrange}\textbf{[statistical functions | Hard]}}\\[2pt]
For \texttt{a = np.array([[1, 2], [3, 4]])}, what does \texttt{a.mean(axis=0)} return?

\begin{enumerate}[label=(\alph*),topsep=2pt,itemsep=1pt,partopsep=0pt,parsep=0pt]
    \item \relax \texttt{array([1.5, 2.5, 3.5])}
    \item \relax \texttt{array([1.5, 3.5])}
    \item \relax \texttt{array([2., 3.])}
    \item \relax \texttt{2.5}
\end{enumerate}
\vspace{2mm}

\noindent\textbf{Q52.} \hfill {\small\color{LPUOrange}\textbf{[statistical functions | Hard]}}\\[2pt]
What does \texttt{np.std(np.array([2, 4, 4, 4, 5, 5, 7, 9]))} return (population standard deviation)?

\begin{enumerate}[label=(\alph*),topsep=2pt,itemsep=1pt,partopsep=0pt,parsep=0pt]
    \item \relax \texttt{4.0}
    \item \relax \texttt{2.14}
    \item \relax \texttt{2.0}
    \item \relax \texttt{5.0}
\end{enumerate}
\vspace{2mm}

\noindent\textbf{Q53.} \hfill {\small\color{LPUOrange}\textbf{[statistical functions | Hard]}}\\[2pt]
Given an array containing a \texttt{np.nan}, which function ignores the NaN when computing the mean?

\begin{enumerate}[label=(\alph*),topsep=2pt,itemsep=1pt,partopsep=0pt,parsep=0pt]
    \item \relax \texttt{np.nanmean}
    \item \relax \texttt{np.mean}
    \item \relax \texttt{np.mean} with \texttt{axis=None}
    \item \relax \texttt{np.average}
\end{enumerate}
\vspace{2mm}

\noindent\textbf{Q54.} \hfill {\small\color{LPUOrange}\textbf{[statistical functions | Hard]}}\\[2pt]
What is the difference between \texttt{a.var()} and \texttt{a.var(ddof=1)} for a sample array?

\begin{enumerate}[label=(\alph*),topsep=2pt,itemsep=1pt,partopsep=0pt,parsep=0pt]
    \item \relax \texttt{ddof=1} divides by N +
    \item \relax \texttt{ddof=1} divides by N -
    \item \relax They are identical for all arrays
    \item \relax \texttt{ddof=1} computes standard deviation instead of variance
\end{enumerate}
\vspace{2mm}

\noindent\textbf{Q55.} \hfill {\small\color{LPUOrange}\textbf{[statistical functions | Hard]}}\\[2pt]
For \texttt{a = np.array([[1, 5], [3, 2]])}, what does \texttt{np.argmax(a)} return?

\begin{enumerate}[label=(\alph*),topsep=2pt,itemsep=1pt,partopsep=0pt,parsep=0pt]
    \item \relax \texttt{3}
    \item \relax \texttt{5}
    \item \relax \texttt{(0, 1)}
    \item \relax \texttt{1}
\end{enumerate}
\vspace{2mm}

\noindent\textbf{Q56.} \hfill {\small\color{LPUOrange}\textbf{[broadcasting | Hard]}}\\[2pt]
Given \texttt{a} of shape \texttt{(3, 1)} and \texttt{b} of shape \texttt{(1, 4)}, what is the shape of \texttt{a + b}?

\begin{enumerate}[label=(\alph*),topsep=2pt,itemsep=1pt,partopsep=0pt,parsep=0pt]
    \item \relax Raises a broadcasting error
    \item \relax \texttt{(3, 1)}
    \item \relax \texttt{(3, 4)}
    \item \relax \texttt{(1, 4)}
\end{enumerate}
\vspace{2mm}

\noindent\textbf{Q57.} \hfill {\small\color{LPUOrange}\textbf{[broadcasting | Hard]}}\\[2pt]
Which pair of shapes CANNOT be broadcast together?

\begin{enumerate}[label=(\alph*),topsep=2pt,itemsep=1pt,partopsep=0pt,parsep=0pt]
    \item \relax \texttt{(5, 3, 4)} and \texttt{(3, 4)}
    \item \relax \texttt{(3, 1)} and \texttt{(1, 4)}
    \item \relax \texttt{(3, 4)} and \texttt{(2, 4)}
    \item \relax \texttt{(3, 4)} and \texttt{(4,)}
\end{enumerate}
\vspace{2mm}

\noindent\textbf{Q58.} \hfill {\small\color{LPUOrange}\textbf{[broadcasting | Hard]}}\\[2pt]
What does the following produce? 

\begin{pycode}
import numpy as np\\
a = np.array([[1], [2], [3]])\\
b = np.array([10, 20, 30])\\
print((a + b).shape)
\end{pycode}

\begin{enumerate}[label=(\alph*),topsep=2pt,itemsep=1pt,partopsep=0pt,parsep=0pt]
    \item \relax \texttt{(1, 3)}
    \item \relax \texttt{(9,)}
    \item \relax \texttt{(3, 3)}
    \item \relax \texttt{(3, 1)}
\end{enumerate}
\vspace{2mm}

\noindent\textbf{Q59.} \hfill {\small\color{LPUOrange}\textbf{[broadcasting | Hard]}}\\[2pt]
To subtract the per-column mean from each element of a \texttt{(100, 5)} array \texttt{X}, which expression works via broadcasting?

\begin{enumerate}[label=(\alph*),topsep=2pt,itemsep=1pt,partopsep=0pt,parsep=0pt]
    \item \relax \texttt{X - X.mean()}
    \item \relax \texttt{X - X.mean(axis=1).reshape(-1, 1)} giving column-mean subtraction
    \item \relax \texttt{X - X.mean(axis=0)}
    \item \relax \texttt{X - X.mean(axis=1)}
\end{enumerate}
\vspace{2mm}

\noindent\textbf{Q60.} \hfill {\small\color{LPUOrange}\textbf{[broadcasting | Hard]}}\\[2pt]
What is the result of the following? 

\begin{pycode}
import numpy as np\\
a = np.array([1, 2, 3])\\
b = np.array([[0], [1], [2]])\\
print(a * b)
\end{pycode}

\begin{enumerate}[label=(\alph*),topsep=2pt,itemsep=1pt,partopsep=0pt,parsep=0pt]
    \item \relax \texttt{[[0 2 6]]}
    \item \relax \texttt{[[1 2 3], [1 2 3], [1 2 3]]}
    \item \relax Raises a broadcasting error
    \item \relax \texttt{[[0 0 0], [1 2 3], [2 4 6]]}
\end{enumerate}
\vspace{2mm}


\newpage
\section{Section B: Comprehensive Subjective Problem with Detailed Analytical Rubric}

\begin{subjectiveq}{4}{1}
Demonstrate NumPy vectorization, multi-dimensional array manipulation, and statistical broadcasting:
\begin{enumerate}[label=(\alph*)]
    \item Given a $1000 \times 10$ sensor dataset, normalize each feature column to zero-mean and unit variance ($z$-score standardization) using vectorized operations without any explicit Python \texttt{for} loops.
    \item Implement min-max feature scaling $[0, 1]$ utilizing column-wise broadcasting.
    \item Compute pairwise Euclidean distances between two matrices $X \in \mathbb{R}^{m \times d}$ and $Y \in \mathbb{R}^{n \times d}$ using NumPy broadcasting.
\end{enumerate}
\end{subjectiveq}

\begin{solbox}
\textbf{Vectorized NumPy Solution:}
\begin{pycode}
import numpy as np\\
~\\
def standardize\_matrix(X):\\
\quad mean = np.mean(X, axis=0) \quad \# shape: (10,)\\
\quad std = np.std(X, axis=0) \quad \quad \# shape: (10,)\\
\quad std\_safe = np.where(std == 0, 1.0, std)\\
\quad return (X - mean) / std\_safe \quad \# Broadcasts (1000, 10) - (10,)\\
~\\
def min\_max\_scale(X):\\
\quad x\_min = np.min(X, axis=0)\\
\quad x\_max = np.max(X, axis=0)\\
\quad diff = np.where(x\_max == x\_min, 1.0, x\_max - x\_min)\\
\quad return (X - x\_min) / diff\\
~\\
def pairwise\_euclidean(X, Y):\\
\quad \# X: (m, 1, d), Y: (1, n, d) -> diff: (m, n, d)\\
\quad diff = X[:, np.newaxis, :] - Y[np.newaxis, :, :]\\
\quad return np.sqrt(np.sum(diff ** 2, axis=-1))
\end{pycode}
\textbf{Broadcasting Invariants:} Alignment occurs along trailing dimensions. Operating in C-contiguous memory blocks yields 50x--100x acceleration over pure Python iterations due to SIMD vector pipelines and cache spatial locality.
\end{solbox}

\begin{examinertip}
\begin{itemize}[leftmargin=1.5em, itemsep=2pt]
    \item \textbf{Algorithmic Formulation (3 Marks):} Correct problem decomposition, function signatures, and parameter documentation.
    \item \textbf{Syntactic Rigor \& Implementation (4 Marks):} Valid Python syntax, proper error handling, defensive input validation, and avoidance of anti-patterns.
    \item \textbf{Complexity \& Edge Analysis (3 Marks):} Accurate time and space asymptotic analysis $\mathcal{O}(n)$ and testing edge conditions (null inputs, zero division, empty collections).
\end{itemize}
\end{examinertip}
